% Chapter 6 content — included by main.tex; compile chapter6.tex for preview.
% Reference style (like IEEE paper): **Label:** regular body text on the same line.
% - Label before colon: bold (bfseries)
% - Following text: normal weight (mdseries / normalfont)
\makeatletter
\@ifundefined{md@runinlabel}{%
  \newcommand{\md@runinlabel}[1]{{\bfseries #1}\normalfont}%
  \newcommand{\mdsubsubsection}[2][]{%
    \refstepcounter{subsubsection}%
    \par\vspace{0.35ex}%
    \noindent\md@runinlabel{\arabic{subsubsection})~#2:}\space
    \def\@mdss@label{#1}%
    \ifx\@mdss@label\@empty\else\label{#1}\fi
  }%
  \newcommand{\mdparagraph}[2][]{%
    \par\vspace{0.25ex}%
    \noindent\md@runinlabel{#2:}\space
    \def\@mdpar@label{#1}%
    \ifx\@mdpar@label\@empty\else\label{#1}\fi
  }%
}{}
\makeatother


\section{Evaluation}
\label{sec:chapter6}

\subsection{Evaluation Overview}

This chapter evaluates the performance of the proposed EPFD-Aware Beam Reassociation for Service Recovery (EABR) framework, focusing on its ability to improve user satisfaction and service continuity under the aggregate EPFD constraint. The evaluation covers five main aspects, which are presented in six subsections. First, we observe the aggregate EPFD variation during the critical period near the GS, as well as the number of beams that must be shut down to satisfy the EPFD constraint. This analysis confirms the existence of EPFD violation and beam shutdown demand during the critical period, and further supports the rationale of designing a reassociation-based service recovery framework. Second, we examine helper satellite availability under different constellation densities to determine whether helper satellite candidates are generally available under the evaluated constellation densities. Third, we compare different interference mitigation methods in terms of service performance and service recovery capability. Fourth, we evaluate the online computational overhead of the proposed framework under different user loads. Finally, we analyze the impact of different EPFD limits on user satisfaction to evaluate the stability and applicability of EABR under different interference constraint conditions.

\subsection{Simulation Setup}

\begin{table}[t]
\centering
\caption{Simulation parameters.}
\label{tab:ch5_sim_params}
\scriptsize
\setlength{\tabcolsep}{3pt}
\begin{tabular}{@{}p{0.54\columnwidth}p{0.38\columnwidth}@{}}
\hline
\textbf{Parameter} & \textbf{Value} \\
\hline
\multicolumn{2}{@{}l@{}}{\textbf{LEO constellation}} \\
\hline
Altitude (km) & 1200 \\
Orbit plane inclination & $87.9^\circ$ \\
Time slot & 1\,s \\
Maximum transmit power & 16.8\,W per satellite \\
Maximum antenna gain & 34.33\,dBi \\
Downlink frequency & 11.7\,GHz \\
Bandwidth & 250\,MHz \\
Antenna pattern & \cite{jalali2023joint} $A\exp(\beta\phi)$ \\
Maximum tilt angle & $10^\circ$ \\
Approximated antenna gain coefficient, ($A$) & $4.0\times 10^{3}$ \\
Approximated antenna gain coefficient, ($\beta$) & $-0.0671$ \\
\hline
\multicolumn{2}{@{}l@{}}{\textbf{GSO satellite}} \\
\hline
Altitude (km) & 35785.4 \\
Longitude & 30.6 \\
Latitude & 0 \\
Downlink frequency & 11.7\,GHz \\
\hline
\multicolumn{2}{@{}l@{}}{\textbf{GSO ground station}} \\
\hline
Longitude & 30.6 \\
Latitude & 0 \\
Minimum elevation angle & $10^\circ$ \\
Maximum receive gain & 38.1\,dBi \\
Noise temperature & 240\,K \\
Antenna pattern & \cite{itur_s1428_2001} \\
\hline
\multicolumn{2}{@{}l@{}}{\textbf{LEO user ground station}} \\
\hline
Maximum receive gain & 40.95\,dBi \\
Noise temperature & 240\,K \\
Antenna pattern & \cite{itur_s1428_2001} \\
User demand & 50\,Mbps \\
\hline
\multicolumn{2}{@{}l@{}}{\textbf{EPFD}} \\
\hline
EPFD limit (dB(W/m$^2$/40\,kHz)) & $-173.4$~\cite{itur_rr2020} \\
Reference antenna diameter & 60\,cm \\
\hline
\end{tabular}
\end{table}

To evaluate the performance of the proposed EABR framework under the NGSO--GSO coexistence scenario, we design a simulation framework that combines STK and MATLAB. This framework is used to construct the LEO constellation, obtain the time-varying satellite geometry, and further evaluate the performance of different interference mitigation methods in terms of EPFD compliance, downlink capacity, and user satisfaction.

The simulation uses the following two main tools:
\begin{itemize}[leftmargin=*,nosep]
  \item \textbf{STK (Satellite Tool Kit):}~\cite{agi2023stk}
  STK is used to construct the OneWeb-like LEO constellation and generate the relative positions, visibility relationships, and beam coverage information among the LEO satellites, the GS, and users at each time slot.
  \item \textbf{MATLAB:}~\cite{mathworks2023matlab}
  MATLAB is used to process the system information generated by STK and perform offline preprocessing. This includes constructing the beam--user association, critical/helper satellite relationships, time-slot beam role records, and candidate beam relation graphs. MATLAB is then used to execute the proposed EABR framework and calculate the EPFD and user satisfaction performance of each method at different time slots.
\end{itemize}

The main simulation parameters used for the OneWeb-like service-performance evaluation are summarized in Table~\ref{tab:ch5_sim_params}. The OneWeb-like constellation and beam configuration are established based on~\cite{oneweb2016fcc}. The table includes the settings related to the LEO constellation, GSO satellite, GS, LEO user ground station, and EPFD evaluation. These parameters are used in the following EPFD evaluation and user satisfaction comparison.

\subsection{Compared Schemes}

To evaluate the proposed EABR framework, we compare it with several interference mitigation methods. All methods are evaluated under the same simulation parameters to ensure a fair comparison. The compared schemes are described as follows:
\begin{itemize}[leftmargin=*,nosep]
  \item \textbf{Beam shutdown only:}
  When the aggregate EPFD exceeds the limit, the high-interference beams on critical satellites are directly shut down to satisfy the EPFD constraint. This method does not perform beam reassociation for service recovery.
  \item \textbf{PC + Tilt:}
  This baseline follows the joint power and satellite tilt control method proposed by Jalali et al.~\cite{jalali2023joint}. The method jointly adjusts the transmit power and tilt angle of LEO satellites to reduce interference toward the GS while satisfying the aggregate EPFD constraint. It is selected as the main comparison method because it considers both EPFD compliance and user service performance through power and attitude adjustment. However, it does not use neighboring helper beams or beam reassociation to recover users affected by interference mitigation.
  \item \textbf{Only HBR:}
  This method only applies HBR. Helper recovery beams serve users located under critical closed beams using the default power of the helper satellite. SBR and EPFD-constrained power allocation are not applied.
  \item \textbf{EABR (Proposed):}
  This method includes SBR, EPFD-constrained power allocation for helper recovery beams, and HBR.
\end{itemize}

\subsection{Evaluation Results}
\label{sec:eval_results}

Based on the simulation settings described above, this section evaluates EABR in terms of EPFD compliance, helper satellite availability, service recovery performance, and online computational overhead. The analysis is divided into six subsections. First, the aggregate EPFD and the number of closed beams are examined to verify whether EABR can still satisfy the EPFD constraint after reassociation-based service recovery, and to observe the beam shutdown behavior of critical satellites at different time slots. Second, we examine different constellation densities to see whether critical satellites usually have helper satellite candidates. Third, the average satisfaction of different methods is compared under different user loads to analyze how increasing system load affects service performance. Fourth, the number of recovered closed-beam users is used to explain the main reason why EABR improves satisfaction, and the satisfaction CDF is further used to observe the satisfaction distribution of all users. Fifth, the online execution times of Only HBR, EABR, and the reproduced PC + Tilt implementation are evaluated under different user loads to examine their computational overhead. Finally, the average user satisfaction of EABR under different EPFD limits is evaluated to observe how the strictness of the interference limit affects service performance.

\mdsubsubsection{EPFD Compliance and Beam Shutdown}

\begin{figure}[t]
  \centering
  \includegraphics[width=1\columnwidth]{ch5_epfd_constraint}
  \caption{Aggregate EPFD over time during the critical period. The y-axis uses a 40\,kHz reference bandwidth, and the EPFD limit is $-173.4$\,dB(W/m$^2$/40\,kHz).}
  \label{fig:ch5_epfd_constraint}
\end{figure}

\begin{figure}[t]
  \centering
  \includegraphics[width=1\columnwidth]{ch5_FullPower_BeamShutdown}
  \caption{Number of closed beams on critical satellites over time during the critical period.}
  \label{fig:ch5_beam_shutdown}
\end{figure}

This subsection first analyzes the aggregate EPFD after applying EABR and further observes the number of beams that must be shut down on critical satellites to satisfy the EPFD constraint. This analysis is used to verify whether the proposed framework can still satisfy the EPFD limit after reassociation-based service recovery, and to show the beam shutdown behavior during the critical period.

Fig.~\ref{fig:ch5_epfd_constraint} shows the aggregate EPFD variation during the critical period near the GS. All EPFD values in this figure use the same 40\,kHz reference bandwidth as Table~\ref{tab:ch5_sim_params}. The red line represents the EPFD constraint, and the black line represents the aggregate EPFD after applying EABR. As shown in the figure, EABR keeps the aggregate EPFD below the limit during the entire observation period. This indicates that the proposed framework does not cause EPFD violation while performing service recovery for closed-beam users. Therefore, EABR can compensate for the service loss caused by beam shutdown while maintaining EPFD compliance.

Fig.~\ref{fig:ch5_beam_shutdown} shows the number of beams that must be shut down on critical satellites at different time offsets. When critical satellites approach the worst EPFD slot, the number of closed beams gradually increases and reaches its maximum around the worst EPFD slot. After the satellites move away from the GS, the number of closed beams gradually decreases. This result shows that the beam shutdown demand changes with the satellite--GS geometry.

\mdsubsubsection{Helper Satellite Availability under Different Constellation Densities}

This subsection investigates how satellite distribution density affects the availability of helper satellite candidates. Three scenarios are considered: a OneWeb-like reference scenario, a high-density scenario, and a low-density scenario. The OneWeb-like scenario serves as the reference case. The satellite distributions of the high-density and low-density scenarios are based on Starlink and Lightspeed, respectively~\cite{starlink2021fcc,telesat2022fcc}. Only their satellite distributions are used as references for setting the different density scenarios.

The GS is located at $(0^\circ\mathrm{N},\,30.6^\circ\mathrm{E})$, consistent with Table~\ref{tab:ch5_sim_params}, and the sub-satellite point of the target GSO satellite is located at the same latitude and longitude as the GS. For each density scenario, the time when one LEO satellite passes directly above the GS is used as the reference point. During the period from 120\,s before to 120\,s after this reference point, the critical satellites and their helper satellite candidates are recorded in each time slot.

\begin{figure*}[t]
  \centering
  \includegraphics[width=\textwidth]{ch6_density_worst_epfd_coverage}
  \caption{Critical and helper satellite coverage at the worst EPFD slot under different satellite distribution densities. (a) Low-density scenario. (b) OneWeb-like reference scenario. (c) High-density scenario. The internal lines represent the 16 satellite-fixed beams.}
  \label{fig:ch6_density_worst_epfd_coverage}
\end{figure*}

Fig.~\ref{fig:ch6_density_worst_epfd_coverage} shows the satellite coverage at the worst EPFD slot under the three density scenarios. The critical satellites and their coverage regions are marked in red, while the helper satellites and their coverage regions are marked in blue. The yellow triangle represents the GS. In the OneWeb-like reference scenario, one critical satellite is surrounded by several helper satellites. In the high-density scenario, multiple critical satellites and a larger number of helper satellites appear around the GS, resulting in more coverage overlap. In the low-density scenario, one critical satellite and two helper satellites appear around the GS at the worst EPFD slot. Although the number of nearby helper satellites is smaller than in the other scenarios, the critical satellite still has helper coverage from both the north and south sides. Since the figure shows only the worst EPFD slot, the two helper satellites shown in the figure may differ from the average value reported in Table~\ref{tab:helper_availability}.

Table~\ref{tab:helper_availability} compares helper satellite availability under different satellite distribution densities using the following four metrics:
\begin{itemize}[leftmargin=*,nosep]
  \item \textbf{Avg. critical satellites per critical slot:}
  The average number of critical satellites in a time slot containing at least one critical satellite.
  \item \textbf{Avg. helper candidates per critical satellite:}
  The average number of helper satellite candidates available to each critical satellite.
  \item \textbf{Critical-satellite helper availability:}
  The percentage of critical satellite instances that have at least one helper satellite candidate.
  \item \textbf{Closed-beam helper-overlap availability:}
  The percentage of closed beams that overlap with at least one helper recovery beam.
\end{itemize}

\begin{table}[t]
\centering
\caption{Helper Satellite Availability under Different Constellation Densities}
\label{tab:helper_availability}
\scriptsize
\setlength{\tabcolsep}{3pt}
\begin{tabular}{@{}lcccc@{}}
\hline
\textbf{\shortstack{Density\\ scenario}} & \textbf{\shortstack{Avg.\\ critical\\ satellites\\ per critical\\ slot}} & \textbf{\shortstack{Avg.\\ helper\\ candidates\\ per critical\\ satellite}} & \textbf{\shortstack{Critical-\\satellite\\ helper\\ availability\\ (\%)}} & \textbf{\shortstack{Closed-\\beam\\ helper-\\overlap\\ avail. (\%)}} \\
\hline
OneWeb-like reference & 1.00 & 4.89 & 100.00 & 100.00 \\
High-density          & 2.97 & 19.03 & 100.00 & 100.00 \\
Low-density           & 1.00 & 1.80 & 100.00 & 69.19  \\
\hline
\end{tabular}
\end{table}

As shown in Table~\ref{tab:helper_availability}, the OneWeb-like reference scenario has an average of 1.00 critical satellite per critical slot, and each critical satellite has an average of 4.89 helper satellite candidates. Both the critical-satellite helper availability and the closed-beam helper-overlap availability reach 100\%.

The high-density scenario has an average of 2.97 critical satellites per critical slot. Although more critical satellites may appear in the same time slot, each critical satellite has an average of 19.03 helper satellite candidates. Both availability metrics also reach 100\%. This result indicates that the denser satellite distribution provides more helper satellite choices around each critical satellite.

In the low-density scenario, an average of 1.00 critical satellite appears in each critical slot, and each critical satellite has an average of 1.80 helper satellite candidates. The critical-satellite helper availability remains at 100\%, indicating that every critical satellite has at least one helper satellite candidate during the evaluated period. However, the closed-beam helper-overlap availability decreases to 69.19\%. This means that although every critical satellite has a helper satellite candidate, some closed beams do not overlap with any helper recovery beam.

Overall, satellite distribution density affects both the number of helper satellite candidates and the closed-beam coverage provided by helper recovery beams. The high-density scenario provides the largest number of helper satellite candidates, while the low-density scenario provides fewer helper choices. The critical-satellite helper availability remains at 100\% in all three scenarios, showing that every critical satellite has at least one helper satellite candidate. However, the lower closed-beam helper-overlap availability in the low-density scenario shows that a helper satellite may not cover all closed beams of a critical satellite. Therefore, a lower satellite distribution density may limit the service recovery coverage of EABR, even when helper satellites are available.

\mdsubsubsection{User Satisfaction under Different User Loads}

\begin{figure*}[t]
\centering
\begin{minipage}[b]{0.32\textwidth}
  \centering
  \includegraphics[width=\linewidth]{ch5_U30_AvgUserSatisfaction_4MethodCompare}
  \centerline{(a)}
\end{minipage}
\hfill
\begin{minipage}[b]{0.32\textwidth}
  \centering
  \includegraphics[width=\linewidth]{ch5_U50_AvgUserSatisfaction_4MethodCompare}
  \centerline{(b)}
\end{minipage}
\hfill
\begin{minipage}[b]{0.32\textwidth}
  \centering
  \includegraphics[width=\linewidth]{ch5_U70_AvgUserSatisfaction_4MethodCompare}
  \centerline{(c)}
\end{minipage}
\caption{Average user satisfaction of different interference mitigation methods under different user loads. (a) 30 users. (b) 50 users. (c) 70 users.}
\label{fig:ch5_avg_satisfaction_loads}
\end{figure*}

This subsection compares the average satisfaction of different interference mitigation methods under different user loads. By changing the number of users, we can observe how the service performance changes as the system load increases, and evaluate whether EABR can maintain user satisfaction under different load conditions.

Fig.~\ref{fig:ch5_avg_satisfaction_loads} shows the average satisfaction under three user load scenarios: 30 users, 50 users, and 70 users. As shown in the figure, Beam shutdown only has the lowest satisfaction because it does not provide service recovery for users affected by closed beams. PC + Tilt can reduce the interference toward the GS. HBR can serve part of the closed-beam users through helper recovery beams, so it achieves better satisfaction than Beam shutdown only and PC + Tilt.

Compared with the other methods, EABR maintains higher average satisfaction under all three user loads. When the user load is low, the helper satellite usually has enough available power, so the difference between Only HBR and EABR is small. However, as the user load increases, more closed-beam users need service recovery, and Only HBR becomes limited by the default power of the helper satellite. In contrast, EABR uses SBR to release available power from the helper satellite and improve the service recovery capability of helper recovery beams. Therefore, EABR can maintain more stable service performance under medium and high user loads.

\mdsubsubsection{Service Recovery Capability and Satisfaction Distribution}

This subsection further analyzes why EABR improves user satisfaction and observes the satisfaction distribution of different methods. Since 30 users is a low-load scenario, the default available power of the helper satellite is usually sufficient to support most recovery demand. Therefore, the difference in service recovery capability is not obvious under 30 users. To better show the effect of the proposed framework, this subsection focuses on the 50-user and 70-user scenarios.

\begin{figure*}[t]
\centering
\begin{minipage}[b]{0.48\textwidth}
  \centering
  \includegraphics[width=\linewidth]{ch5_U50_RecoveredUserCount}
  \centerline{(a)}
\end{minipage}
\hfill
\begin{minipage}[b]{0.48\textwidth}
  \centering
  \includegraphics[width=\linewidth]{ch5_U70_RecoveredUserCount}
  \centerline{(b)}
\end{minipage}
\caption{Number of recovered closed-beam users achieved by Only HBR and EABR under different user loads. (a) 50 users. (b) 70 users.}
\label{fig:ch5_recovered_user_count}
\end{figure*}

Fig.~\ref{fig:ch5_recovered_user_count} compares the number of successfully recovered closed-beam users between Only HBR and EABR under the 50-user and 70-user scenarios. As shown in the figure, EABR serves more closed-beam users in both load scenarios. The difference becomes more obvious around the worst EPFD slot, where the beam shutdown demand is the most severe. This result shows the necessity of SBR in EABR.

\begin{figure*}[t]
\centering
\begin{minipage}[b]{0.48\textwidth}
  \centering
  \includegraphics[width=\linewidth]{ch5_U50_UserSatisfaction_CDF_3MethodCompare}
  \centerline{(a)}
\end{minipage}
\hfill
\begin{minipage}[b]{0.48\textwidth}
  \centering
  \includegraphics[width=\linewidth]{ch5_U70_UserSatisfaction_CDF_3MethodCompare}
  \centerline{(b)}
\end{minipage}
\caption{CDF of user satisfaction for different interference mitigation methods under different user loads. (a) 50 users. (b) 70 users.}
\label{fig:ch5_satisfaction_cdf}
\end{figure*}

Fig.~\ref{fig:ch5_satisfaction_cdf} shows the satisfaction CDF of three interference mitigation methods. This figure focuses on the user-level satisfaction distribution. A CDF curve that shifts further to the right means that more users maintain higher satisfaction. In contrast, if the curve rises quickly in the low-satisfaction region, more users suffer from service degradation.

As shown in the figure, the CDF curve of EABR shifts more toward the high-satisfaction region than the other methods. This means that more users can maintain better service quality. The result shows that EABR not only improves average satisfaction, but also reduces the proportion of low-satisfaction users, making the overall user service quality more stable.

\mdsubsubsection{On-Orbit Computational Overhead}

This subsection evaluates the on-orbit computational overhead of Only HBR, EABR, and the reproduced PC + Tilt implementation under different user loads. The per-slot procedure is updated once in each 1\,s time slot. A 1\,s time interval has also been used in previous NGSO--GSO interference simulations~\cite{itur_s1655_2003}. Therefore, the 1\,s time-slot duration is used as the reference for evaluating the per-slot execution time. The evaluation focuses on the period from 30\,s before to 30\,s after the worst EPFD slot. Since the number of closed beams varies across time slots, this period includes different levels of service recovery demand. Only HBR, EABR, and the reproduced PC + Tilt implementation are executed separately. Their average, minimum, and maximum execution times are recorded under each user-load scenario.

\begin{figure}[t]
  \centering
  \includegraphics[width=1\columnwidth]{ch5_overhead_only_hbr_vs_eabr}
  \caption{Per-slot execution time of Only HBR and EABR under different user loads. Each bar shows the average execution time per time slot, and the vertical black lines indicate the minimum and maximum execution times over the evaluation period.}
  \label{fig:ch5_overhead_only_hbr_vs_eabr}
\end{figure}

As shown in Fig.~\ref{fig:ch5_overhead_only_hbr_vs_eabr}, Only HBR has a lower execution time under all user loads because it only performs the reassociation of closed-beam users. In contrast, EABR additionally performs SBR and EPFD-constrained power allocation, resulting in higher online computational overhead. As the number of users increases, the average execution time of EABR increases from 141.42~ms for 30 users to 299.78~ms for 50 users and 513.63~ms for 70 users. This trend indicates that higher user loads and a larger number of reassociation decisions increase the computational cost of EABR. For the maximum execution time, EABR requires 214.44~ms and 461.94~ms under the 30-user and 50-user scenarios, respectively, both of which are below the 1~s time-slot duration adopted in this work. When the number of users increases to 70, the maximum execution time reaches 1066.11~ms, slightly exceeding 1~s in the worst observed time slot. Overall, the average execution time of EABR remains below 1~s under all evaluated user loads.

\begin{figure}[t]
  \centering
  \includegraphics[width=1\columnwidth]{runtime_overhead_pc_tilt_vs_eabr}
  \caption{Per-slot execution time of EABR and the reproduced PC + Tilt implementation under different user loads. Each bar shows the average execution time per time slot, and the vertical black lines indicate the minimum and maximum execution times over the evaluation period.}
  \label{fig:ch6_overhead_pc_tilt_vs_eabr}
\end{figure}

Fig.~\ref{fig:ch6_overhead_pc_tilt_vs_eabr} further compares the per-slot execution time of EABR and our reproduced implementation of PC + Tilt. Since the original PC + Tilt study does not report its execution time, the results shown in this figure are obtained from our reproduced implementation. As shown in Fig.~\ref{fig:ch6_overhead_pc_tilt_vs_eabr}, the reproduced PC + Tilt implementation has a higher average execution time than EABR under all evaluated user loads. It also shows a larger execution-time range. This result indicates that EABR has lower online computational overhead under the evaluated settings.

Overall, Only HBR has the lowest online computational overhead because it only performs closed-beam user reassociation. EABR requires more computation because it also performs SBR and EPFD-constrained power allocation, but its average execution time remains below the 1\,s time-slot duration under all evaluated user loads. Compared with our reproduced implementation of PC + Tilt, EABR requires less execution time and shows a smaller execution-time range. Therefore, under the simulation settings of this work, EABR provides a better balance between service recovery capability and online computational overhead.

\mdsubsubsection{Impact of Different EPFD Limits}

This subsection further evaluates the impact of different EPFD limits on the service performance of EABR. When the EPFD limit becomes stricter, the allowed aggregate interference becomes lower. As a result, critical satellites may need to shut down more beams to satisfy the constraint, and the recovery demand of closed-beam users may also increase. This analysis is used to observe whether EABR can maintain stable average user satisfaction under different EPFD limits.

\begin{figure}[t]
  \centering
  \includegraphics[width=1\columnwidth]{ch5_U70_EABR_AvgUserSatisfaction_EpfdCompare}
  \caption{Average user satisfaction of EABR under different EPFD limits. The compared EPFD limits are all defined with the same 40\,kHz reference bandwidth as Table~\ref{tab:ch5_sim_params}.}
  \label{fig:ch5_epfd_limit_compare}
\end{figure}

Fig.~\ref{fig:ch5_epfd_limit_compare} shows the average user satisfaction of EABR under different EPFD limits. All compared limits use the same 40\,kHz reference bandwidth. As shown in the figure, when the EPFD limit is more relaxed, EABR can maintain higher average user satisfaction. This result shows that the strictness of the EPFD limit directly affects user satisfaction. It also indicates that EABR can achieve better service performance under more relaxed interference constraints.

Overall, EABR keeps the aggregate EPFD below the limit during the critical period. The helper availability analysis under different constellation densities further shows that helper satellite candidates are generally available under the evaluated constellation densities. EABR then uses SBR to improve the ability of helper recovery beams to serve closed-beam users. Compared with other interference mitigation methods, EABR maintains higher average satisfaction under different user loads and reduces the proportion of low-satisfaction users. The online computational overhead results further show that the average execution time of EABR remains below the 1-s time-slot duration under all evaluated user loads. In addition, the results under different EPFD limits show that the strictness of the interference constraint affects user satisfaction, and EABR can maintain better service performance under more relaxed constraints. These results show that EABR achieves a better balance among EPFD compliance, service continuity, and online computational overhead.
